\chapter{EG-CNTFET stability}
\label{cap:chapter3}

\input{./chapters/chapter3/stability.tex}

\newpage
\thispagestyle{empty}
\ % Empty page
\newpage

%%%%%%%%%%%%%%%%%%%%%%%%%%%%%%%%%%%%%%%%%%%%%%%%%%%%%%%%%%%%%%%%%%%%%%

\section{Standard EG-FET}
\label{sec:big_channel}

\input{./chapters/chapter3/bigChannel}

%%%%%%%%%%%%%%%%%%%%%%%%%%%%%%%%%%%%%%%%%%%%%%%%%%%%%%%%%%%%%%%%%%%%%%

\section{Reduced gate-channel distance}
\label{sec:small_distance}

\input{./chapters/chapter3/smallDistance}

%%%%%%%%%%%%%%%%%%%%%%%%%%%%%%%%%%%%%%%%%%%%%%%%%%%%%%%%%%%%%%%%%%%%%%

\section{Reduced channel area}
\label{sec:small_channel}

\input{./chapters/chapter3/smallChannel}

%%%%%%%%%%%%%%%%%%%%%%%%%%%%%%%%%%%%%%%%%%%%%%%%%%%%%%%%%%%%%%%%%%%%%%

\section{Conclusion}
\label{sec:stability_conclusion}

In conclusion, this chapter explored the properties of three EG-CNTFET geometries to understand their stabilization mechanisms. 

It has been observed that for the standard platform with exposed carbon
nanotubes, the current diminishes rapidly; however, it also has the best \vth{}
and \ratio{} values. If a lipophilic membrane is added on the channel to protect
the CNT network, stabilization improves significantly, as it is achieved in just
30 minutes, and is maintained for up to 12 hours of testing. On the other hand,
there are no positive effects when the semiconductor channel is moved closer to
the gate, and performance actually deteriorates when the device size is reduced.
A more thorough analysis of the results obtained indicates that gating is
primarily governed by capacitive EDL effects, with small gate currents
indicating proper EDL formation even in membrane-encapsulated devices.
Hysteresis can most probably be attributed to charge trapping and detrapping
phenomena, when mobile ions and charges move under bias before finding a
steady-state. The presence of a lipophilic membrane significantly alters device
characteristics by modifying ionic coupling and doping levels, resulting in
\vth{} shifts and reduced \ratio{}. Specifically, a lower ionic coupling reduces
EDL capacitance, leading to lower transconductance and thus lower drain current.
Additionally, the membrane increases drain-source resistance, therefore lowering
\ion{}, which in turn lowers the \ratio{}. However, the membrane provides better
stability by preventing CNT doping by oxygen and water; this also slows down ion
permeation (stabilization phase), which eventually occurs, leading to a more
stable device. Finally, scaling effects were tested, by studying the behaviour
of the smaller devices: these exhibited inferior performance due to the lower
conductivity of the channel and to the higher impact of small defects on overall
device behavior.




%%%%%%%%%%%%%%%%%%%%%%%%%%%%%%%%%%%%%%%%%%%%%%%%%%%%%%%%%%%%%%%%%%%%%%

\newpage
\thispagestyle{empty}
\ % Empty page
\newpage
